\section{A Combinatorial Bound}

The probabilistic analysis needed for our later results can be abstracted into a single
lemma. The proof of the lemma is independent of the rest of the paper.

We define the function $L(h,S)$ for any real number $h$ greater than one and any positive
integer $S$ as follows: Let $L(h,S)$ be the smallest integer such that in $L(h,S)$ independent
Bernoulli trials each with probability $h$ of success, the probability of having fewer than $S$
successes is less than $h^{-1}$.

The following simple upper bound holds for the whole range of values of $S$ and $h$ and shows
that $L(h,S)$ is essentially linear both in $h$ and in $S$.

\medskip

\noindent\textbf{PROPOSITION:}\quad For \emph{all} integers $S \geqslant 1$ and \emph{all} real $h > 1$,
\[
  L(h,S) \leqslant 2h\,(S + \log_e h) \;.
\]

\noindent\textit{Proof.}\quad We use the following three well known inequalities of which the last is due to Chernoff
(see~\cite{5}, p.~18).

\begin{enumerate}
\item[(a)] For all \quad $x > 0$ \quad $(1 + x^{-1})^x < e$ \;.
\item[(b)] For all \quad $x > 0$ \quad $(1 - x^{-1})^x < e^{-1}$ \;.
\item[(c)] In $m$ independent trials each with probability $p$ of success the probability that
there are at most $k$ successes, where $k < mp$, is at most
\[
  \binom{m - mp}{m - k}^{m-k} \binom{mp}{k}^k .
\]
\end{enumerate}

The first factor in the expression in (c) above can be rewritten as
\[
  \left(1 - \frac{mp - k}{m - k}\right)^{((m-k)/(mp-k))\,(mp-k)} ,
\]
using (b) with $x = (m{-}k)/(mp{-}k)$ we can upper bound the product by
\[
  e^{-mp+k}\,(mp/k)^k \;.
\]
Substituting $m = 2h(S + \log_e h)$, $p = h^{-1}$ and $k = S$ and using logarithms to the base $e$ gives the bound
\[
  e^{-2S - 2\log h + S} \cdot \bigl[2\,(1 + (\log h)/S)\bigr]^{(S/\log h)\,\log h} \;.
\]
Rewriting this using (a) with $x \approx (\log_e h)/S$ gives
\[
  e^{-S - 2\log h} \cdot 2^S \cdot e^{\log h} \leqslant (2/e)^S \cdot e^{-\log h} \leqslant h^{-1} . \qquad \qed
\]

As an illustration of an application suppose that we have access to EXAMPLES, a source of
natural positive examples of vectors for concept $F$. Suppose that we want to find an approximation
to the subset $P$ of variables that are determined in at least one natural example of $F$. Consider
the following procedure. Pick the first $L(h, |P|)$ vectors provided by EXAMPLES and let $P'$ be the
union of the sets of variables determined by these vectors. The proposition then implies that with
probability at least $1 - h^{-1}$ the set $P'$ will have the following property: if a vector is
chosen with probability distribution $D$ then the probability that it determines a variable in
$P - P'$ is less than $4h^{-1}$. To see this observe that if $P'$ does not have the claimed property then the
following has happened: $L(h, |P|)$ trials have been made (i.e.\ calls of EXAMPLES) each with probability
greater than $h^{-1}$ of success (i.e.\ finding a vector that determines a variable in $P - P'$) but
there have been fewer than $|P|$ successes (i.e.\ discoveries of new additions to $P'$). The probability of such a happening is less than $h^{-1}$ by the definition of $L$.

The above application shows that the set of variables that are determined in natural examples
of $F$ can be approximated by a procedure whose \textbf{runtime} is independent of $t$, the total number of
variables.
